% GENERATED FILE. Edit canonical lessons, then run npm run generate:books.
% canonical-source-hash: fnv1a64:b0cba9ebee497583
% canonical-lessons: TE-C33-anuko, TE-C33-artham-cesuko, TE-C33-caduvu, TE-C33-raayu

\chapter{Four Verbs of the Mind}
\label{ch:te-verbs-of-the-mind}
\hlchaptermodality{33}

\begin{chapteropening}
\textbf{By the end of this chapter:} \emph{I can say I think, I understand, I read and I write in Telugu with the same three slots, build the -\te{కో} ending that hands an action back to its doer, say \textquotedblleft{}I understood\textquotedblright{} with no \te{నేను} in it, and tell this chapter's native verbs from the Sanskrit words layered beside them.}

Say \te{నేను} \te{తెలుగు} \te{రాస్తాను}, run the chapter's four verbs in the present, count them up from \te{ఆరు} to \te{పది}, take \te{వ్రాయు} back to the Dravidian line-drawing root behind Tamil \ta{வரை} and Kannada \ka{ಬರೆ}, and name which of the four words came from Sanskrit.
\end{chapteropening}

\section[anukō]{\te{అనుకో} (anukō) — \textquotedblleft{}to think,\textquotedblright{} which is saying it to yourself}

\label{lesson:TE-C33-anuko}



\begin{quote}
\emph{Nāku telusu} — \textquotedblleft{}I know\textquotedblright{} — and that verb had no room for a person. This one takes the person back, and carries a second verb on its end.
\end{quote}

\subsection*{You'll want to know: \te{అనుకో}}
\begin{quote}
\textbf{\te{అనుకో}} — \emph{anukō} — \textbf{to think; to suppose}
\end{quote}

Threaded up: \textbf{\te{అనుకుంటాను}} (\emph{anukuṇṭānu}), \textquotedblleft{}I think,\textquotedblright{} and \textbf{\te{అనుకున్నాను}} (\emph{anukunnānu}), \textquotedblleft{}I thought.\textquotedblright{} That \textbf{-\te{ంటాను}} is the ending you already made in \textbf{\te{తింటాను}} (\emph{tiṇṭānu}): stem, tense, then the \emph{-ānu} that means \textquotedblleft{}I.\textquotedblright{}

\subsection*{The letters in this word}
\textbf{\te{అ}} is the standing vowel \emph{a}, used when a word begins on that sound. \textbf{\te{ను}} is \emph{na} with the \emph{u} sign, the shape that ends \emph{nēnu}. \textbf{\te{కో}} is \textbf{\te{క}} (\emph{ka}) with the long-\emph{ō} sign: hold it, \emph{kō}.

\begin{grammarlens}[title={Grammar Lens: -\te{కో}, the ending that hands the verb back}]
\textbf{\te{కొను}} (\emph{konu}) is a verb of its own — \textquotedblleft{}to take, to get.\textquotedblright{} Bolted onto the back of another verb it stops being a word and becomes an ending, \textbf{-\te{కొను}}, worn down to \textbf{-\te{కో}} in speech. What it adds is small and everywhere: \emph{do the thing, and keep it}.

So \textbf{\te{అను}} (\textquotedblleft{}say\textquotedblright{}) plus \textbf{-\te{కో}} is \textbf{say it, and keep it} — say it to nobody but yourself.
\end{grammarlens}

\begin{cousinweb}
\textbf{\te{అను}} (\emph{anu}), \textquotedblleft{}to say,\textquotedblright{} is Proto-Dravidian \emph{\textbf{*aHn-}}, and every sister kept it: Tamil \textbf{\ta{என்}} (\emph{eṉ}), Kannada \textbf{\ka{ಅನ್ನು}} (\emph{annu}), Malayalam \textbf{\ml{എന്നുക}} (\emph{ennuka}).

Beside it Telugu keeps a borrowed thinker: \textbf{\te{ఆలోచించు}} (\emph{ālōcin̄cu}), \textquotedblleft{}to consider\textquotedblright{} — \textbf{\te{ఆలోచిస్తాను}} (\emph{ālōcistānu}). Its Sanskrit noun \textbf{\te{ఆలోచన}} (\emph{ālōcana}) is glossed two ways at once: \textquotedblleft{}reflecting\textquotedblright{} \textbf{and} \textquotedblleft{}seeing.\textquotedblright{} Sanskrit made thinking a kind of looking, as English did in \emph{reflect}.

The split Chapter 32 named over \emph{tiṇḍi} and \emph{bhōjanaṁ} holds: plain verb native, formal word Sanskrit.
\end{cousinweb}

\subsection*{Your turn}
Say these aloud:

\begin{itemize}
\raggedright
  \item the two pieces, then the word — \textquotedblleft{}anu … kō … anukō\textquotedblright{}
  \item \textquotedblleft{}anukuṇṭānu\textquotedblright{}, I think — then \textquotedblleft{}anukunnānu\textquotedblright{}, I thought
  \item the ending you already had — \textquotedblleft{}tiṇṭānu … anukuṇṭānu\textquotedblright{}
  \item the family's say-verb — \textquotedblleft{}anu … eṉ … annu … ennuka\textquotedblright{}
  \item native and borrowed — \textquotedblleft{}anukuṇṭānu … ālōcistānu\textquotedblright{}
\end{itemize}

\begin{tcolorbox}[breakable,colback=teal!4,colframe=teal!35!black,title={Before you move on}]
What are the two pieces of \emph{anukō}, and what does the second do? (\emph{\textbf{\te{అను}}} \textquotedblleft{}say\textquotedblright{} plus \textbf{-\te{కో}}, which hands the action back to the doer.) Say \textquotedblleft{}I think,\textquotedblright{} and name the verb ending the same way. (\emph{Anukuṇṭānu}; \textbf{\te{తిను}}.) Which sisters keep the say-verb, and what is \emph{ālōcana} glossed as? (\textbf{All three}; \textquotedblleft{}\textbf{reflecting}\textquotedblright{} and \textquotedblleft{}\textbf{seeing}\textquotedblright{}.) Which think-word is borrowed? (\textbf{\te{ఆలోచించు}} — as \emph{bhōjanaṁ} was beside \emph{tiṇḍi}.)
\end{tcolorbox}
\section[arthaṁ cēsukō]{\te{అర్థం} \te{చేసుకో} (arthaṁ cēsukō) — \textquotedblleft{}to understand,\textquotedblright{} making the meaning yours}

\label{lesson:TE-C33-artham-cesuko}



\begin{quote}
\emph{Anukuṇṭānu} — I think — and the \textbf{-\te{కో}} on its back meant \emph{and keep it}. That ending now arrives on a verb you have owned since Chapter 5.
\end{quote}

\subsection*{You'll want to know: \te{అర్థం} \te{చేసుకో}}
\begin{quote}
\textbf{\te{అర్థం} \te{చేసుకో}} — \emph{arthaṁ cēsukō} — \textbf{to understand}
\end{quote}

Three pieces, all already yours: \textbf{\te{అర్థం}} (\emph{arthaṁ}, \textquotedblleft{}meaning\textquotedblright{}), \textbf{\te{చేయు}} (\emph{cēyu}, \textquotedblleft{}to do\textquotedblright{}), and the \textbf{-\te{కో}} of the last lesson. Word for word it says \textbf{make the meaning your own}.

\begin{quote}
\textbf{\te{నేను} \te{అర్థం} \te{చేసుకుంటాను}.} — \emph{nēnu arthaṁ cēsukuṇṭānu} — \textquotedblleft{}I understand.\textquotedblright{}
\end{quote}

\subsection*{The letters in this word}
\textbf{\te{ర్థ}} is a \textbf{conjunct} — \textbf{\te{ర}} (\emph{ra}) stacked above \textbf{\te{థ}} (\emph{tha}), one block for two consonants with no vowel between. \textbf{\te{ం}} is the \textbf{anusvāra} that closes \emph{namaskāram}, and \textbf{\te{చే}} opens \emph{cēyu}.

\begin{grammarlens}[title={Grammar Lens: the other way to say it, with no \textquotedblleft{}I\textquotedblright{} in it}]
\begin{quote}
\textbf{\te{నాకు} \te{అర్థమైంది}.} — \emph{nāku arthamaindi} — \textquotedblleft{}I understood.\textquotedblright{} Literally: \textquotedblleft{}\textbf{to me it became meaning}.\textquotedblright{}
\end{quote}

Three dative sentences now, and one design:

\noindent
\begin{tabularx}{\linewidth}{@{}>{\raggedright\arraybackslash}X>{\raggedright\arraybackslash}X@{}}
\toprule
\textbf{the sentence} & \textbf{what it says} \\
\midrule
\emph{nāku telugu vaccu} & Telugu \textbf{comes} to me \\
\emph{nāku telusu} & it \textbf{is clear} to me \\
\emph{nāku arthamaindi} & it \textbf{became meaning} to me \\
\bottomrule
\end{tabularx}

No \emph{nēnu} in any of them: the person sits in the dative because understanding, like knowing, arrives. Both shapes are good Telugu and do not feel alike — \emph{nēnu arthaṁ cēsukuṇṭānu} is work you did, \emph{nāku arthamaindi} is a thing that landed.
\end{grammarlens}

\begin{cousinweb}
\textbf{\te{అర్థం}} is Sanskrit \textbf{\dv{अर्थ}} (\emph{artha}), and Telugu took the whole of it: \textbf{\te{అర్థము}} means \textquotedblleft{}\textbf{meaning}\textquotedblright{} and \textquotedblleft{}\textbf{money, wealth}\textquotedblright{} in one entry — which is why Kauṭilya's manual of statecraft is the \emph{Arthaśāstra}. Its opposite: \textbf{\te{అపార్థం}} (\emph{apārthaṁ}), \textquotedblleft{}a meaning taken wrong.\textquotedblright{}

Now the seam. Chapter 20's \textbf{\te{వాతావరణం}} was Sanskrit throughout — \emph{vāta} plus \emph{āvaraṇa}. This is not that. Only the noun crossed over; the doing is Telugu's own \emph{cēyu} wearing Telugu's own \emph{-kō}.
\end{cousinweb}

\subsection*{Your turn}
Say these aloud:

\begin{itemize}
\raggedright
  \item the three pieces — \textquotedblleft{}arthaṁ … cēyu … kō\textquotedblright{}
  \item \textquotedblleft{}nēnu arthaṁ cēsukuṇṭānu\textquotedblright{} — I understand
  \item the three dative sentences — \textquotedblleft{}nāku telugu vaccu … nāku telusu … nāku arthamaindi\textquotedblright{}
  \item the two verbs built with -kō — \textquotedblleft{}anukuṇṭānu … arthaṁ cēsukuṇṭānu\textquotedblright{}
  \item \textquotedblleft{}arthaṁ … apārthaṁ\textquotedblright{}
\end{itemize}

\begin{tcolorbox}[breakable,colback=teal!4,colframe=teal!35!black,title={Before you move on}]
Take \emph{arthaṁ cēsukō} apart and give the literal sense. (\textbf{Meaning + do + for yourself}.) Say \textquotedblleft{}I understood\textquotedblright{} without \emph{nēnu}, then name the other two sentences built that way. (\emph{Nāku arthamaindi}; \emph{nāku telugu vaccu}, \emph{nāku telusu}.) Why is the person in the dative? (\textbf{Understanding arrives at you}.) What two things does \emph{arthaṁ} mean, and how is this compound unlike \emph{vātāvaraṇam}? (\textbf{Meaning} and \textbf{wealth}; \emph{vātāvaraṇam} is \textbf{Sanskrit throughout}, here only the noun is.)
\end{tcolorbox}
\section[caduvu]{\te{చదువు} (caduvu) — \textquotedblleft{}to read,\textquotedblright{} which is also \textquotedblleft{}an education\textquotedblright{}}

\label{lesson:TE-C33-caduvu}



\begin{quote}
\emph{Nāku arthamaindi} — it became meaning to me. Say it once. Now the verb for the work that gets you there.
\end{quote}

\subsection*{You'll want to know: \te{చదువు}}
\begin{quote}
\textbf{\te{చదువు}} — \emph{caduvu} — \textbf{to read; to study}
\end{quote}

\begin{quote}
\textbf{\te{నేను} \te{తెలుగు} \te{చదువుతాను}.} — \emph{nēnu telugu caduvutānu} — \textquotedblleft{}I read Telugu.\textquotedblright{}
\end{quote}

Telugu does not split reading from studying, and it builds no separate noun either: \textbf{\te{చదువు}} unchanged is \textquotedblleft{}\textbf{education, studies}\textquotedblright{} — what a family means when it says a child's \emph{caduvu} is going well.

The sisters scatter here as completely as they did over seeing: Tamil \textbf{\ta{படி}} (\emph{paḍi}), Kannada \textbf{\ka{ಓದು}} (\emph{ōdu}), Malayalam \textbf{\ml{വായിക്കുക}} (\emph{vāyikkuka}), Telugu \textbf{\te{చదువు}}. Four unrelated everyday verbs for one act.

\subsection*{The letters in this word}
\textbf{\te{చ}} is \emph{ca}, the letter that opens \emph{cūḍu}. \textbf{\te{దు}} is \textbf{\te{ద}} (\emph{da}) with the \emph{u} sign, and \textbf{\te{వు}} is \textbf{\te{వ}} (\emph{va}) with the same sign. Three light syllables, each ending in that short \emph{u}.

\begin{grammarlens}[title={Grammar Lens: one stem, three middle vowels}]
\noindent
\begin{tabularx}{\linewidth}{@{}>{\raggedright\arraybackslash}X>{\raggedright\arraybackslash}X@{}}
\toprule
\textbf{when} & \textbf{Telugu} \\
\midrule
now, or later & \emph{caduvutānu} \\
already & \emph{cadivānu} \\
asked politely & \emph{cadavaṇḍi} \\
\bottomrule
\end{tabularx}

\emph{cadu-}, \emph{cadi-}, \emph{cada-}. The middle vowel moves, and nothing else does.

Chapter 32 gave you the harsh version of this with \textbf{\te{రా}}: there the word you call the verb by was not the word the endings attach to, and they went to \emph{vaccu-} instead. This is the mild version — one stem, three vowels, plainly the same word. Expect the wobble; do not expect a second verb.

The respectful \textbf{-\te{అండి}} arrives exactly as it did on \emph{cūḍaṇḍi}: \textbf{\te{చదవండి}} (\emph{cadavaṇḍi}), \textquotedblleft{}please read.\textquotedblright{} The ending never changes; only the stem in front of it does.
\end{grammarlens}

\subsection*{Your turn}
Say these aloud:

\begin{itemize}
\raggedright
  \item \textquotedblleft{}nēnu telugu caduvutānu\textquotedblright{} — I read Telugu
  \item the three middle vowels — \textquotedblleft{}caduvutānu … cadivānu … cadavaṇḍi\textquotedblright{}
  \item the two polite commands — \textquotedblleft{}cūḍaṇḍi … cadavaṇḍi\textquotedblright{}
  \item the four sisters' four verbs — \textquotedblleft{}caduvu … paḍi … ōdu … vāyikkuka\textquotedblright{}
  \item the chapter so far — \textquotedblleft{}anukuṇṭānu, arthaṁ cēsukuṇṭānu, caduvutānu\textquotedblright{}
\end{itemize}

\begin{tcolorbox}[breakable,colback=teal!4,colframe=teal!35!black,title={Before you move on}]
What two English verbs does \emph{caduvu} cover, and what is its noun? (\textbf{Read} and \textbf{study}; the \textbf{same word} — an education.) Say \textquotedblleft{}I read Telugu,\textquotedblright{} then \textquotedblleft{}please read.\textquotedblright{} (\emph{Nēnu telugu caduvutānu}; \emph{cadavaṇḍi}.) What moves between \emph{caduvutānu}, \emph{cadivānu} and \emph{cadavaṇḍi}, and which verb warned you about a second stem? (\textbf{The middle vowel}; \textbf{\te{రా}}, whose endings went to \emph{vaccu-}.) Do the four sisters share a read-verb? (\textbf{No} — \emph{caduvu}, \emph{paḍi}, \emph{ōdu}, \emph{vāyikkuka}.)
\end{tcolorbox}
\section[rāyu]{\te{రాయు} (rāyu) — \textquotedblleft{}to write,\textquotedblright{} which began as drawing a line}

\label{lesson:TE-C33-raayu}



\begin{quote}
\emph{Caduvutānu}, \emph{cadivānu}, \emph{cadavaṇḍi} — one stem, three middle vowels. The chapter's last verb is the other half of that one.
\end{quote}

\subsection*{You'll want to know: \te{రాయు}}
\begin{quote}
\textbf{\te{రాయు}} — \emph{rāyu} — \textbf{to write}
\end{quote}

\begin{quote}
\textbf{\te{నేను} \te{తెలుగు} \te{రాస్తాను}.} — \emph{nēnu telugu rāstānu} — \textquotedblleft{}I write Telugu.\textquotedblright{}
\end{quote}

The past is \textbf{\te{రాశాను}} (\emph{rāśānu}). Dictionaries often print the headword as \textbf{\te{వ్రాయు}} (\emph{vrāyu}): not a different verb, just the older spelling, with a \emph{v} gone quiet in speech.

\subsection*{The letters in this word}
\textbf{\te{రా}} is \textbf{\te{ర}} (\emph{ra}) — Chapter 32's whole command \emph{rā} — with the long-\emph{ā} sign, and \textbf{\te{యు}} is \textbf{\te{య}} (\emph{ya}) with the \emph{u} sign. In \textbf{\te{వ్రాయు}} a \textbf{\te{వ}} is stacked in front.

\begin{cousinweb}
That silent \emph{v} is the evidence. \textbf{\te{వ్రాయు}} goes back to a Dravidian \emph{\textbf{*warV-}}, \textquotedblleft{}to draw a line,\textquotedblright{} and the sisters kept the \emph{v} Telugu buried: Tamil \textbf{\ta{வரை}} (\emph{varai}), Malayalam \textbf{\ml{വരയ്ക്കുക}} (\emph{varaykkuka}), Kannada \textbf{\ka{ಬರೆ}} (\emph{bare}) with a \emph{b}. Every one still means \textbf{draw} as readily as \textbf{write}.

Telugu's noun keeps the shape: \textbf{\te{రాత}} (\emph{rāta}), \textquotedblleft{}handwriting\textquotedblright{} — and, in speech, what is written for you.

One honest note: Tamil's everyday write-verb is \textbf{\ta{எழுது}} (\emph{eḻutu}), not \emph{varai}. The sisters share the root without agreeing which word does the daily work — the picture \emph{caduvu} left you with.
\end{cousinweb}

\begin{grammarlens}[title={Grammar Lens: ten verbs, and one machine}]
Chapter 32 handed you \textbf{\te{ఆరు}} (\emph{āru}) verbs. Four more makes \textbf{\te{పది}} (\emph{padi}), ten — and not one needed a new machine. Stem, tense, person: the three slots were already yours. That is what agglutination buys.

What they added is a shape and a habit. The shape is \textbf{-\te{కో}}, which hands the action back to its doer. The habit is Telugu's oldest: \emph{caduvu}, \emph{rāyu} and \emph{anu} are native, while \emph{ālōcana} and \emph{arthaṁ}, laid on top, are Sanskrit — \emph{tiṇḍi} beside \emph{bhōjanaṁ}, one chapter on.
\end{grammarlens}

\subsection*{Your turn}
Say these aloud:

\begin{itemize}
\raggedright
  \item \textquotedblleft{}nēnu telugu rāstānu\textquotedblright{} — then \textquotedblleft{}rāśānu\textquotedblright{}
  \item the spellings and the sisters — \textquotedblleft{}rāyu … vrāyu … varai … bare\textquotedblright{}
  \item the chapter's four — \textquotedblleft{}anukuṇṭānu, arthaṁ cēsukuṇṭānu, caduvutānu, rāstānu\textquotedblright{}
  \item the count — \textquotedblleft{}āru\textquotedblright{}, then \textquotedblleft{}padi\textquotedblright{}
  \item two dative sentences — \textquotedblleft{}nāku telusu … nāku arthamaindi\textquotedblright{}
\end{itemize}

\begin{tcolorbox}[breakable,colback=teal!4,colframe=teal!35!black,title={Before you move on}]
Say \textquotedblleft{}I write Telugu,\textquotedblright{} name the older spelling, and say what the root meant first. (\emph{Nēnu telugu rāstānu}; \textbf{\te{వ్రాయు}}; \textbf{to draw a line} — Tamil \emph{varai}, Kannada \emph{bare}.) Run this chapter's four verbs, then count them with Chapter 32's. (\emph{Anukuṇṭānu, arthaṁ cēsukuṇṭānu, caduvutānu, rāstānu}; \textbf{padi} — \emph{āru} plus four.) What does \textbf{-\te{కో}} do, and which two carry it? (\textbf{Hands the action back to the doer}; \emph{anukō}, \emph{arthaṁ cēsukō}.) Which words here are Sanskrit? (\emph{\textbf{\te{అర్థం}}} and \emph{ālōcana}, with the plain verbs native — \emph{tiṇḍi} against \emph{bhōjanaṁ}.)
\end{tcolorbox}
